\cvsection{Expériences professionnelles}

\begin{cventries}
  \cventry
    {Stagiaire ingénieur de recherche en IA}
    {INRIA}
    {Bordeaux, France}
    {Mars 2026 -- Septembre 2026}
    {
      \begin{cvitems}
        \item {Analyse d’un besoin des chercheurs et développement en autonomie d’une application fonctionnelle intégrant un agent IA multi-étapes pour la recherche de publications scientifiques.}
        \item {Conception d’un pipeline RAG pour traiter et interroger un corpus documentaire, avec chunking, embeddings, indexation dans ChromaDB et recherche sémantique.}
        \item {Développement d’un backend Python/FastAPI avec API REST documentées via OpenAPI et stockage des documents et métadonnées dans MongoDB.}
        \item {Création d’une interface React/Next.js permettant aux chercheurs d’interroger l’agent et de consulter les documents retrouvés.}
        \item {Développement de l’agent sans framework d’orchestration, avec intégration de l’API OpenAI, conception de prompts adaptés au cas d’usage et gestion des appels d’outils et du contexte.}
        \item {Automatisation de la collecte de métadonnées de publications en ligne avec des tâches cron et transmission des résultats par webhooks.}
        \item {Définition d’un protocole expérimental et comparaison de plusieurs LLM et quantifications sur cluster HPC selon la réussite de la recherche et la latence.}
        \item {Conteneurisation de composants avec Docker et mise en place d’une CI/CD sur GitHub.}
      \end{cvitems}
    }

  \cventry
    {Stagiaire ingénieur de recherche en apprentissage par renforcement}
    {Université Karunya}
    {Coimbatore, Inde}
    {Juin 2025 -- Septembre 2025}
    {
      \begin{cvitems}
        \item {Développement d’un modèle et de son environnement d’entraînement pour la navigation autonome d’un robot hospitalier, puis déploiement et tests sur Raspberry Pi.}
      \end{cvitems}
    }
\end{cventries}
